% !PRINTED 469
\weberpage{469}

und es sind also acht numerische Co\"efficienten zu berechnen, die
sich durch Vergleichung einiger Glieder in den Ausdr\"ucken der
$a_\nu$ durch die Wurzeln einerseits, durch die Invarianten (\S.\,123)
andererseits finden lassen, und die ausser rationalen Zahlen nur
$\sqrt{-7}$ enthalten k\"onnen.

Wir wollen diese Co\"efficienten zun\"achst nur unter der
Voraussetzung berechnen, dass $f=0$ ist ${}^1$).

Man braucht dann nur die ersten Glieder (mit $x_1^5x_2$, $x_1^{10}x_2^2$)
in den Potenzsummen $\Sigma z_r^3$ und $\Sigma z_r^6$ zu berechnen und die
\gesperrt{Newton}'schen Formeln anzuwenden. Der letzte Co\"efficient ergiebt
sich aus dem einem Gliede $x_1^{14}$ in dem Producte der $z_r$. Man
findet zun\"achst
\[
\Sigma z_r^3 = -3\cdot7\cdot\frac{1-\sqrt{-7}}{2}\,x_1^5x_2\dots
\]
\[
\Sigma z_r^6 = \left[6\cdot7 + 15\cdot7\left(\frac{1-\sqrt{-7}}{2}\right)^2\right]x_1^{10}x_2^2\dots
\]
und daraus die gesuchte Resolvente
\weq{9}{z^7 - 7\cdot\frac{1-\sqrt{-7}}{2}\,\Delta z^4 - 7\cdot\frac{5+\sqrt{-7}}{2}\,\Delta^2 z - C = 0.}
Will man diese Gleichung auf eine andere zur\"uckf\"uhren, die
nur von den Verh\"altnissen der $x_1, x_2, x_3$ abh\"angt, so setzt man
\weq{10}{z = \frac{Cu}{\Delta^2}, \quad g = \frac{\Delta^7}{C^3},}
wodurch man aus (9) eine Gleichung erh\"alt, in der nur noch der
eine Parameter $g$ vorkommt, n\"amlich
\weq{11}{u^7 - 7\,\frac{1-\sqrt{-7}}{2}\,g\,u^4 - 7\,\frac{5+\sqrt{-7}}{2}\,g^2\,u - g^2 = 0.}
Macht man dieselbe Substitution in der allgemeinen Resol\-vente,
in der $f$ nicht $=0$ gesetzt ist, so hat man noch einen
weiteren Parameter
\weq{12}{h = \frac{f\Delta^4}{C^2}}
einzuf\"uhren, und die Co\"efficienten der Resolvente werden, von
%CONT>>

\par\noindent\rule{0.3\linewidth}{0.4pt}\par
${}^1$) Dieser Fall ist darum von besonderem Interesse, weil er, \"ahnlich
wie die Ikosa\"edergleichung, auf die Transformationsgleichungen aus der
Theorie der elliptischen Functionen f\"uhrt.
